% =====================================================================
%  Banco complementar --- Campo Magnetico  (REVISADO)
%  Substitui a versao gerada por template (5 clones em N1, 9 em N2,
%  8 em N3 e 8 questoes conceituais marcadas como N4).
%  O cap11 ja cobre: B proporcional a 1/r (MC), fio de 10 A a 5 cm,
%  regra da mao direita, solenoide de 1000 esp/m, ESPIRA CIRCULAR de
%  R=10 cm no centro, fio de 4 A a 2 cm, DOIS FIOS PARALELOS (mesmo
%  sentido e sentidos opostos), TOROIDE e duas espiras concentricas.
%  Este banco evita esses casos e explora: fio de secao finita por
%  lei de Ampere, cabo coaxial, campo no eixo de espira, dipolo
%  magnetico a grande distancia, projeto de solenoide, mapeamento com
%  sensor Hall, blindagem magnetica e os limites praticos de eletroimas.
%  Distribuicao mantida: 5 N1 + 9 N2 + 8 N3 + 8 N4 = 30
% =====================================================================

\subsubsection*{Banco complementar --- Campo magnético}

\begin{atencao}[Organização do banco]
Três resultados cobrem quase tudo: fio retilíneo longo
($B=\mu_0I/2\pi r$), centro de espira ($B=\mu_0I/2R$) e solenoide longo
($B=\mu_0nI$, \emph{independente} do raio). Note que apenas o primeiro decai com
a distância --- dentro do solenoide o campo é uniforme. O N3 trata de condutores
de seção finita, cabo coaxial e campo fora do eixo; o N4 exige dedução por
Biot--Savart, projeto, blindagem e limites físicos.
\end{atencao}

\blocofixacao

\begin{questao}[1]{}
Um fio retilíneo longo conduz $\SI{2}{\ampere}$. Determine o campo magnético a
$\SI{1}{\centi\meter}$ dele.
\dados{$\mu_0=4\pi\times10^{-7}\,\si{\tesla\meter\per\ampere}$}
\resposta{$B=\SI{40}{\micro\tesla}$}
\resolucao{
\[
B=\frac{\mu_0 I}{2\pi r}=\frac{(4\pi\times10^{-7})(2)}{2\pi(0{,}01)}
=\frac{\pot{2}{-7}\times2}{0{,}01}=\pot{4}{-5}\,\si{\tesla}.
\]
\textbf{Atalho útil:} $\dfrac{\mu_0}{2\pi}=\pot{2}{-7}$ exatamente, de modo que
$B=\pot{2}{-7}\dfrac{I}{r}$ --- evita carregar $\pi$ pela conta toda.\\
\textbf{Ordem de grandeza:} $\SI{40}{\micro\tesla}$ é comparável ao campo
magnético da Terra ($\approx\SI{50}{\micro\tesla}$). Uma bússola perto desse fio
já seria perturbada.}
\end{questao}

\begin{questao}[1]{}
Um solenoide tem $\SI{2000}{}$ espiras por metro e conduz
$\SI{3}{\ampere}$. Determine o campo em seu interior.
\resposta{$B=\SI{7.54}{\milli\tesla}$}
\resolucao{
\[
B=\mu_0 nI=(4\pi\times10^{-7})(2000)(3)=\pot{7.54}{-3}\,\si{\tesla}.
\]
\textbf{Note o que \emph{não} aparece na fórmula:} o raio do solenoide e o
número total de espiras. Só importa a \textbf{densidade linear}
$n=N/L$.\\
Dois solenoides, um fino e um grosso, com a mesma densidade de espiras e a mesma
corrente, produzem o \emph{mesmo} campo interno.}
\end{questao}

\begin{questao}[1]{}
Determine a corrente necessária para que um fio retilíneo produza
$B=\SI{100}{\micro\tesla}$ a $\SI{2}{\centi\meter}$.
\resposta{$I=\SI{10}{\ampere}$}
\resolucao{
\[
I=\frac{2\pi rB}{\mu_0}=\frac{rB}{\pot{2}{-7}}
=\frac{(0{,}02)(\pot{1}{-4})}{\pot{2}{-7}}=\SI{10}{\ampere}.
\]
\textbf{Leitura prática:} são necessários $\SI{10}{\ampere}$ para produzir, a
dois centímetros, apenas o dobro do campo terrestre. Campos magnéticos
intensos exigem correntes altas --- ou muitas espiras.}
\end{questao}

\begin{questao}[1]{}
Sobre o campo magnético de um solenoide longo, é correto afirmar:
\begin{alternativas}[2]
\item é uniforme no interior e praticamente nulo no exterior
\item é maior no exterior
\item aumenta com a distância ao eixo, no interior
\item é nulo no interior
\end{alternativas}
\resposta{(a)}
\resolucao{Aplicando a lei de Ampère a um retângulo com um lado dentro e outro
fora, obtém-se $B_{int}=\mu_0nI$ e $B_{ext}\approx0$ para um solenoide
\emph{ideal} (infinitamente longo).\\
\textbf{Por que o campo externo quase se anula:} as linhas que saem de uma
extremidade retornam pela outra, espalhando-se por um volume enorme --- a
densidade de linhas fora é desprezível diante da interna.\\
\textbf{Em solenoides reais}, o campo externo não é exatamente nulo, e nas
extremidades o campo interno cai para cerca de \emph{metade} do valor central.}
\end{questao}

\begin{questao}[1]{}
O campo magnético da Terra vale aproximadamente $\SI{50}{\micro\tesla}$.
Expresse-o em gauss.
\dados{$\SI{1}{\tesla}=10^{4}\,\si{\gauss}$}
\resposta{$\SI{0.5}{\gauss}$}
\resolucao{
\[
B=\pot{5}{-5}\,\si{\tesla}\times10^{4}=\SI{0.5}{\gauss}.
\]
\textbf{Escala de referência para o tópico:}
\begin{center}
\begin{tabular}{lcc}
\hline
Situação & Tesla & Gauss\\
\hline
Campo terrestre & $\pot{5}{-5}$ & $0{,}5$\\
Ímã de geladeira & $\pot{5}{-3}$ & $50$\\
Ímã de neodímio (superfície) & $0{,}5$ & $5000$\\
Ressonância magnética & $1{,}5$ a $3$ & $\pot{1.5}{4}$ a $\pot{3}{4}$\\
\hline
\end{tabular}
\end{center}
O tesla é uma unidade \emph{grande}: campos de laboratório raramente passam de
alguns tesla, e o gauss continua em uso justamente por ser mais próximo dos
valores cotidianos.}
\end{questao}

\blocoaplicacao

\begin{questao}[2]{}
Um solenoide de $100$ espiras e $\SI{10}{\centi\meter}$ de comprimento conduz
$\SI{1}{\ampere}$, com núcleo de ar.
\begin{itensq}
\item Determine a densidade de espiras.
\item Determine o campo interno.
\end{itensq}
\resposta{$n=\SI{1000}{\per\meter}$; $B=\SI{1.26}{\milli\tesla}$}
\resolucao{(a) $n=\dfrac{100}{0{,}10}=\SI{1000}{\per\meter}$.\\
(b) $B=\mu_0nI=(4\pi\times10^{-7})(1000)(1)
=\pot{1.257}{-3}\,\si{\tesla}$.\\
\textbf{Comparação instrutiva:} isso é apenas $25$ vezes o campo terrestre ---
um solenoide modesto produz campo modesto. Para chegar a
$\SI{0.1}{\tesla}$ seriam necessários $\SI{80}{\ampere}$ com a mesma
geometria, o que exigiria refrigeração.\\
\textbf{Caminho alternativo:} inserir núcleo ferromagnético, que multiplica
$B$ por $\mu_r$ --- centenas ou milhares de vezes.}
\end{questao}

\begin{questao}[2]{}
Um fio conduz corrente constante. Compare o campo a $\SI{4}{\centi\meter}$ com
o campo a $\SI{12}{\centi\meter}$.
\resposta{$B_1/B_2=3$}
\resolucao{Como $B\propto1/r$:
\[
\frac{B_1}{B_2}=\frac{r_2}{r_1}=\frac{12}{4}=3.
\]
\textbf{Contraste importante com o campo elétrico:} o campo de uma carga
puntiforme cai com $1/r^{2}$, mas o de um fio \emph{infinito} cai com
$1/r$.\\
A razão é a mesma da eletrostática: o fio é uma fonte \emph{estendida} em uma
dimensão, e a superfície amperiana (um cilindro) tem área proporcional a $r$,
não a $r^{2}$.}
\end{questao}

\begin{questao}[2]{}
Uma espira circular de raio $\SI{5}{\centi\meter}$ deve produzir
$\SI{100}{\micro\tesla}$ em seu centro. Determine a corrente.
\resposta{$I\approx\SI{7.96}{\ampere}$}
\resolucao{Do campo no centro de uma espira, $B=\dfrac{\mu_0I}{2R}$:
\[
I=\frac{2RB}{\mu_0}=\frac{2(0{,}05)(\pot{1}{-4})}{4\pi\times10^{-7}}
=\SI{7.96}{\ampere}.
\]
\textbf{Comparação com o fio reto:} um fio retilíneo a
$\SI{5}{\centi\meter}$ precisaria de $\SI{25}{\ampere}$ para o mesmo campo. A
espira é mais eficiente porque \emph{todos} os seus elementos contribuem no
mesmo sentido no centro --- a razão entre as duas fórmulas é exatamente
$\pi$.}
\end{questao}

\begin{questao}[2]{}
Um solenoide tem $N$ espiras e comprimento $L$. Determine o que ocorre com o
campo interno se $N$ e $L$ forem \textbf{ambos dobrados}.
\resposta{O campo não se altera}
\resolucao{Como $B=\mu_0\dfrac{N}{L}I$, dobrar $N$ e $L$ mantém a razão
$N/L$ inalterada:
\[
B'=\mu_0\frac{2N}{2L}I=B.
\]
\textbf{Interpretação:} o solenoide dobrado equivale a dois solenoides
originais colocados em sequência. Cada trecho "enxerga" a mesma vizinhança de
espiras, e o campo local não muda.\\
\textbf{O que muda:} a região de campo uniforme fica mais longa, e a resistência
do enrolamento dobra --- exigindo o dobro da tensão para a mesma corrente.}
\end{questao}

\begin{questao}[2]{}
Dois fios longos e paralelos, distantes $\SI{8}{\centi\meter}$, conduzem
$\SI{4}{\ampere}$ cada, no \textbf{mesmo} sentido. Determine o campo no ponto
médio entre eles.
\resposta{$B=0$}
\resolucao{Cada fio produz, no ponto médio,
\[
B=\pot{2}{-7}\frac{4}{0{,}04}=\pot{2}{-5}\,\si{\tesla}.
\]
Pela regra da mão direita, com as correntes no \emph{mesmo} sentido, os dois
campos têm sentidos \textbf{opostos} no ponto médio --- e módulos iguais.
Cancelam-se:
\[
B_{total}=0.
\]
\textbf{Cuidado com a intuição:} correntes de mesmo sentido \emph{atraem-se},
mas seus campos se \emph{cancelam} entre elas. Força e campo são coisas
diferentes.\\
\textbf{Fora do par}, ao contrário, os campos se somam --- e é por isso que dois
fios juntos conduzindo no mesmo sentido produzem, de longe, o campo de um único
fio de corrente $2I$.}
\end{questao}

\begin{questao}[2]{}
Um fio reto é imerso em um meio de permeabilidade relativa $\mu_r=200$.
Determine o efeito sobre o campo, comparado ao ar.
\resposta{O campo fica $200$ vezes maior}
\resolucao{Em meio linear,
\[
B=\frac{\mu_r\mu_0I}{2\pi r},
\]
de modo que $B$ é multiplicado por $\mu_r=200$.\\
\textbf{Mecanismo:} o material se \emph{magnetiza}, e seus dipolos alinhados
somam sua própria contribuição ao campo aplicado. Diferentemente do dielétrico
--- que \emph{reduz} o campo elétrico --- o material ferromagnético o
\textbf{amplifica}.\\
\textbf{Ressalva séria:} $\mu_r$ não é constante. Ele cai drasticamente quando o
material satura, e a relação $B\times H$ deixa de ser linear. A fórmula acima
vale apenas em campos baixos.}
\end{questao}

\begin{questao}[2]{}
Determine a distância a que um fio de $\SI{50}{\ampere}$ produz campo igual ao
terrestre ($\SI{50}{\micro\tesla}$).
\resposta{$r=\SI{20}{\centi\meter}$}
\resolucao{
\[
r=\frac{\pot{2}{-7}I}{B}=\frac{(\pot{2}{-7})(50)}{\pot{5}{-5}}
=\SI{0.20}{\meter}.
\]
\textbf{Consequência prática:} a menos de $\SI{20}{\centi\meter}$ de um
condutor de $\SI{50}{\ampere}$, uma bússola aponta mais para o fio do que para
o norte.\\
É por isso que instrumentos magnéticos sensíveis --- e monitores de tubo, na
época em que existiam --- precisam ficar afastados de cabos de potência.}
\end{questao}

\begin{questao}[2]{}
Um toroide e um solenoide têm a mesma densidade de espiras e a mesma corrente.
Compare seus campos internos e externos.
\resposta{Campo interno praticamente igual; o toroide tem campo externo
\emph{rigorosamente} nulo}
\resolucao{\textbf{Campo interno:} ambos valem $\approx\mu_0nI$. No toroide,
$n=N/(2\pi r)$ varia ligeiramente com o raio, de modo que o campo não é
perfeitamente uniforme na seção --- mas para toroides finos a diferença é
pequena.\\
\textbf{Campo externo:} aqui está a diferença essencial.
\begin{itemize}
\item \textbf{Solenoide:} as linhas precisam retornar por fora, e o campo
externo, embora fraco, \emph{existe}.
\item \textbf{Toroide:} as linhas fecham-se \emph{inteiramente} dentro do
núcleo. Aplicando a lei de Ampère a um círculo externo, a corrente líquida
envolvida é \textbf{zero} --- e o campo também.
\end{itemize}
\textbf{Consequência de projeto:} indutores toroidais praticamente não irradiam
nem captam interferência, e podem ser montados lado a lado sem acoplamento
mútuo. É por isso que dominam aplicações de filtragem em fontes chaveadas.}
\end{questao}

\begin{questao}[2]{}
Dois fios perpendiculares entre si, ambos no mesmo plano da folha, conduzem
$\SI{6}{\ampere}$ e $\SI{8}{\ampere}$. Determine o campo resultante em um
ponto a $\SI{3}{\centi\meter}$ do primeiro e $\SI{4}{\centi\meter}$ do
segundo.
\resposta{$B=\SI{56.6}{\micro\tesla}$}
\resolucao{Cada fio produz campo perpendicular ao plano definido por ele e pelo
ponto. Como os fios são perpendiculares entre si e coplanares, ambos os campos
saem (ou entram) \emph{perpendicularmente à folha} --- portanto são
\textbf{paralelos} e se somam algebricamente:
\[
B_1=\pot{2}{-7}\frac{6}{0{,}03}=\SI{40}{\micro\tesla};
\qquad
B_2=\pot{2}{-7}\frac{8}{0{,}04}=\SI{40}{\micro\tesla}.
\]
Se os sentidos coincidirem: $B=\SI{80}{\micro\tesla}$; se forem opostos,
$B=0$.\\
\textbf{Atenção ao caso geral:} se os fios \emph{não} fossem coplanares, os
campos teriam direções diferentes e a soma seria vetorial --- com
$\SI{40}{\micro\tesla}$ perpendiculares, o resultado seria
$40\sqrt2=\SI{56.6}{\micro\tesla}$. \textbf{A geometria decide se a soma é
algébrica ou vetorial}, e é ela que precisa ser estabelecida antes de qualquer
conta.}
\end{questao}

\blocoaprofundamento

\begin{questao}[3]{}
Um condutor cilíndrico maciço de raio $a=\SI{2}{\milli\meter}$ conduz
$\SI{100}{\ampere}$ uniformemente distribuídos na seção.
\begin{itensq}
\item Determine o campo a $\SI{1}{\milli\meter}$, na superfície e a
      $\SI{5}{\milli\meter}$ do eixo.
\item Identifique onde o campo é máximo.
\end{itensq}
\resposta{$5$, $10$ e $\SI{4}{\milli\tesla}$; máximo na superfície}
\resolucao{\textbf{Dentro ($r<a$):} a superfície amperiana envolve apenas a
fração $(r/a)^{2}$ da corrente:
\[
B=\frac{\mu_0 I r}{2\pi a^{2}}
=\pot{2}{-7}\frac{100\,r}{(0{,}002)^{2}}.
\]
Em $r=\SI{1}{\milli\meter}$: $B=\SI{5}{\milli\tesla}$.\\
\textbf{Fora ($r>a$):} toda a corrente é envolvida, e vale a fórmula do fio
fino:
\[
B=\pot{2}{-7}\frac{100}{r}.
\]
Em $r=\SI{5}{\milli\meter}$: $B=\SI{4}{\milli\tesla}$.\\
\textbf{Na superfície ($r=a$):} as duas expressões coincidem em
$B=\SI{10}{\milli\tesla}$ --- o campo é contínuo.\\
(b) \textbf{O máximo está na superfície.} Dentro, $B$ cresce linearmente com
$r$; fora, decai com $1/r$. O pico ocorre exatamente em $r=a$.\\
\textbf{Analogia com a esfera isolante carregada:} lá o campo elétrico também
cresce linearmente dentro e cai com $1/r^{2}$ fora, com máximo na superfície. A
estrutura matemática é a mesma --- muda apenas a lei de Gauss ou de Ampère
aplicada.}
\end{questao}

\begin{questao}[3]{}
Um cabo coaxial tem condutor central de raio $a$ conduzindo $I$ e blindagem
externa conduzindo $I$ em sentido \textbf{oposto}, com $I=\SI{5}{\ampere}$.
\begin{itensq}
\item Determine o campo entre os condutores, a $\SI{2}{\milli\meter}$ do eixo.
\item Determine o campo fora do cabo.
\item Explique a consequência prática.
\end{itensq}
\resposta{$B=\SI{0.5}{\milli\tesla}$ entre os condutores; $B=0$ fora}
\resolucao{(a) Entre os condutores, a superfície amperiana envolve apenas a
corrente central:
\[
B=\pot{2}{-7}\frac{5}{0{,}002}=\pot{5}{-4}\,\si{\tesla}=\SI{0.5}{\milli\tesla}.
\]
(b) \textbf{Fora do cabo}, a superfície amperiana envolve $+I$ e $-I$ --- a
corrente líquida é \textbf{zero}:
\[
\oint\vec B\cdot\mathrm{d}\vec\ell=\mu_0(0)=0
\;\Longrightarrow\;
B=0.
\]
(c) \textbf{Consequência prática: o cabo coaxial não irradia nem capta campo
magnético.} Todo o campo fica confinado entre os condutores.\\
Isso o torna a escolha padrão para transmitir sinais em ambientes com
interferência --- e explica por que um cabo paralelo comum, cujos condutores
estão afastados, é muito pior: ali as correntes opostas não se cancelam
perfeitamente e o par forma uma espira que irradia.\\
\textbf{Mesma ideia, outra forma:} o par trançado obtém efeito semelhante ao
inverter a orientação da espira a cada torção, de modo que as contribuições
sucessivas se anulem.}
\end{questao}

\begin{questao}[3]{}
Uma espira de raio $R=\SI{10}{\centi\meter}$ conduz $\SI{5}{\ampere}$.
\begin{itensq}
\item Determine o campo no centro.
\item Determine o campo no eixo, a $\SI{10}{\centi\meter}$ e a
      $\SI{20}{\centi\meter}$ do centro.
\item Comente a taxa de queda.
\end{itensq}
\resposta{$31{,}4$; $11{,}1$ e $\SI{2.81}{\micro\tesla}$}
\resolucao{(a) $B_0=\dfrac{\mu_0I}{2R}=\dfrac{(4\pi\times10^{-7})(5)}{0{,}20}
=\pot{3.14}{-5}\,\si{\tesla}$.\\
(b) No eixo,
\[
B=\frac{\mu_0IR^{2}}{2\left(R^{2}+x^{2}\right)^{3/2}}.
\]
Em $x=R=\SI{0.10}{\meter}$:
$B=\pot{1.11}{-5}\,\si{\tesla}$, ou seja
$B_0/2^{3/2}=0{,}354\,B_0$.\\
Em $x=2R=\SI{0.20}{\meter}$: $B=\pot{2.81}{-6}\,\si{\tesla}$, cerca de
$0{,}089\,B_0$.\\
(c) \textbf{A queda é rápida.} A apenas um raio de distância, o campo já caiu
para $35\%$; a dois raios, para $9\%$.\\
Para $x\gg R$, a expressão tende a
$B\to\dfrac{\mu_0IR^{2}}{2x^{3}}$ --- decaimento com $1/x^{3}$, típico de
\textbf{dipolo}. É o mesmo expoente do dipolo elétrico, e pela mesma razão: a
espira não tem "monopolo magnético" para produzir termo $1/x^{2}$.}
\end{questao}

\begin{questao}[3]{}
Compare o campo no centro de um solenoide \emph{finito} com o do solenoide
ideal, e determine o campo em sua extremidade.
\begin{itensq}
\item Indique a condição para a aproximação ideal.
\item Determine o campo na extremidade.
\item Discuta a consequência para projeto.
\end{itensq}
\resposta{Na extremidade, $B\approx\mu_0nI/2$ --- metade do valor central}
\resolucao{(a) \textbf{Condição:} $L\gg D$, isto é, comprimento muito maior que
o diâmetro. Na prática, $L/D>10$ garante erro inferior a $1\%$ no centro.\\
(b) \textbf{Na extremidade}, apenas "metade" do solenoide contribui --- as
espiras estão todas de um lado. O resultado exato para solenoide longo é
\[
B_{extremidade}=\frac{\mu_0nI}{2},
\]
exatamente \textbf{metade} do valor central.\\
(c) \textbf{Consequências de projeto.}
\begin{itemize}
\item A região de campo verdadeiramente uniforme é \emph{menor} que o
solenoide: ela se restringe à parte central, tipicamente à metade do
comprimento.
\item Para uniformidade melhor que $1\%$ em uma região de comprimento
$\ell$, o solenoide precisa ter comprimento bem maior que $\ell$ --- ou usar
enrolamento com densidade \emph{corrigida} nas extremidades.
\item A alternativa clássica é o par de \textbf{bobinas de Helmholtz}: duas
bobinas iguais separadas por uma distância igual ao raio, arranjo que anula a
segunda derivada do campo no centro e produz região uniforme com pouco
material.
\end{itemize}}
\end{questao}

\begin{questao}[3]{}
Três fios longos e paralelos, coplanares e igualmente espaçados de
$\SI{5}{\centi\meter}$, conduzem $\SI{10}{\ampere}$ cada, todos no mesmo
sentido. Determine o campo no fio central.
\resposta{$B=0$ no fio central}
\resolucao{O fio central não produz campo sobre si mesmo. Restam as
contribuições dos dois vizinhos, ambos a $\SI{5}{\centi\meter}$:
\[
B_1=B_2=\pot{2}{-7}\frac{10}{0{,}05}=\SI{40}{\micro\tesla}.
\]
Pela regra da mão direita, com os dois vizinhos conduzindo no mesmo sentido e
estando em \emph{lados opostos} do fio central, seus campos ali têm sentidos
\textbf{contrários}:
\[
B=40-40=0.
\]
\textbf{Mas a força não é nula --- ela também é.} Como $\vec F=I\vec\ell\times
\vec B$ e $\vec B=\vec 0$ no fio central, a força sobre ele é zero: as atrações
dos dois vizinhos se equilibram.\\
\textbf{Nos fios externos, porém}, a situação é diferente: cada um sofre atração
resultante para o centro. O conjunto tende a se \emph{comprimir} --- efeito real
em barramentos de alta corrente, que precisam de espaçadores mecânicos.}
\end{questao}

\begin{questao}[3]{}
Um eletroímã deve produzir $\SI{0.1}{\tesla}$ com núcleo de ar, usando
solenoide de $\SI{20}{\centi\meter}$.
\begin{itensq}
\item Determine o produto $NI$ necessário.
\item Avalie a viabilidade com $N=1000$ espiras.
\item Proponha alternativa.
\end{itensq}
\resposta{$NI=\SI{15.9}{\kilo\ampere}$-espira; exigiria $\SI{15.9}{\ampere}$
com $1000$ espiras}
\resolucao{(a) De $B=\mu_0NI/L$:
\[
NI=\frac{BL}{\mu_0}=\frac{(0{,}1)(0{,}20)}{4\pi\times10^{-7}}
=\pot{1.59}{4}\ \text{A-espira}.
\]
(b) Com $N=1000$: $I=\SI{15.9}{\ampere}$.\\
\textbf{Viável, mas problemático.} Mil espiras conduzindo
$\SI{16}{\ampere}$ exigem fio grosso (a densidade de corrente admissível limita
a seção), o que aumenta o volume do enrolamento. E a dissipação é considerável:
com resistência de apenas $\SI{1}{\ohm}$, seriam
$\SI{256}{\watt}$ --- exigindo refrigeração.\\
(c) \textbf{Alternativa: núcleo ferromagnético.} Com $\mu_r=1000$, a mesma
corrente produziria $\SI{100}{\tesla}$ --- valor impossível, porque o material
\textbf{satura} em torno de $1{,}5$ a $\SI{2}{\tesla}$.\\
\textbf{O que o núcleo realmente entrega:} os $\SI{0.1}{\tesla}$ pedidos com
corrente cerca de mil vezes menor ($\SI{16}{\milli\ampere}$), enquanto o
material permanecer na região linear.\\
\textbf{Limite físico:} acima da saturação, o núcleo deixa de ajudar e o
eletroímã volta a se comportar como se fosse de ar. É por isso que campos acima
de $\SI{2}{\tesla}$ exigem bobinas supercondutoras, sem núcleo.}
\end{questao}

\begin{questao}[3]{}
Um cabo de par paralelo tem condutores separados por $\SI{2}{\milli\meter}$,
conduzindo $\SI{1}{\ampere}$ em sentidos opostos.
\begin{itensq}
\item Determine o campo no ponto médio entre eles.
\item Determine o campo a $\SI{10}{\centi\meter}$ do par.
\item Compare com um fio único.
\end{itensq}
\resposta{$\SI{400}{\micro\tesla}$ no meio; $\approx\SI{0.04}{\micro\tesla}$ a
$\SI{10}{\centi\meter}$}
\resolucao{(a) No ponto médio, cada fio dista $\SI{1}{\milli\meter}$, e com
correntes \emph{opostas} os campos se \textbf{somam}:
\[
B=2\times\pot{2}{-7}\frac{1}{0{,}001}=\SI{400}{\micro\tesla}.
\]
(b) A $\SI{10}{\centi\meter}$, os dois fios estão praticamente à mesma
distância e seus campos quase se cancelam. O resultado é aproximadamente o de um
\textbf{dipolo}:
\[
B\approx\frac{\mu_0 I d}{2\pi r^{2}}
=\pot{2}{-7}\frac{(1)(0{,}002)}{(0{,}10)^{2}}
=\pot{4}{-8}\,\si{\tesla}.
\]
(c) \textbf{Um fio único} de $\SI{1}{\ampere}$ produziria, a
$\SI{10}{\centi\meter}$:
$B=\pot{2}{-6}\,\si{\tesla}$ --- \textbf{cinquenta vezes mais}.\\
\textbf{Conclusão:} manter os condutores de ida e volta \emph{juntos} reduz
drasticamente o campo irradiado, porque o par se comporta como dipolo
($1/r^{2}$) em vez de fio isolado ($1/r$). Quanto menor a separação $d$, melhor
--- e é exatamente por isso que se recomenda não separar os condutores de um
mesmo circuito.}
\end{questao}

\begin{questao}[3]{}
Um sensor Hall indica $\SI{2.5}{\milli\volt}$ quando exposto a
$\SI{100}{\milli\tesla}$, com resposta linear e tensão de repouso nula.
\begin{itensq}
\item Determine a sensibilidade.
\item Determine o campo correspondente a uma leitura de
      $\SI{0.6}{\milli\volt}$.
\item Discuta as fontes de erro.
\end{itensq}
\resposta{$\SI{25}{\micro\volt\per\milli\tesla}$; $B=\SI{24}{\milli\tesla}$}
\resolucao{(a)
\[
S=\frac{2{,}5\,\si{\milli\volt}}{100\,\si{\milli\tesla}}
=\SI{25}{\micro\volt\per\milli\tesla}.
\]
(b) $B=\dfrac{0{,}6}{0{,}025}=\SI{24}{\milli\tesla}$.\\
(c) \textbf{Fontes de erro, em ordem de importância.}
\begin{itemize}
\item \textbf{Tensão de \emph{offset}:} sensores reais indicam algo mesmo com
$B=0$, por assimetria construtiva. Valores de dezenas de microvolts são comuns
--- comparáveis ao sinal de alguns militesla. \emph{Correção:} medir com o
sensor blindado e subtrair.
\item \textbf{Deriva térmica:} tanto o \emph{offset} quanto a sensibilidade
variam com a temperatura.
\item \textbf{Orientação:} o sensor responde apenas à componente
\emph{perpendicular} à sua face. Um erro angular $\theta$ introduz fator
$\cos\theta$ --- desprezível para ângulos pequenos, mas $13\%$ a
$\ang{30}$.
\item \textbf{Campo terrestre} somado à medida, se não for compensado.
\end{itemize}
\textbf{Procedimento robusto:} medir com o sensor em uma orientação, depois
girá-lo $\ang{180}$ e medir de novo. A média das duas leituras cancela o
\emph{offset}, e a semidiferença dá o campo.}
\end{questao}

\blocodesafio

\begin{questao}[4]{}
\textbf{(Dedução por Biot--Savart.)} Deduza o campo magnético no centro de uma
espira circular de raio $R$ percorrida por corrente $I$.
\begin{itensq}
\item Aplique a lei de Biot--Savart a um elemento de corrente.
\item Integre e obtenha o resultado.
\item Explique por que a integração é simples neste caso.
\end{itensq}
\resposta{$B=\dfrac{\mu_0 I}{2R}$}
\resolucao{(a) A lei de Biot--Savart estabelece que um elemento
$I\,\mathrm{d}\vec\ell$ produz, a uma distância $\vec r$,
\[
\mathrm{d}\vec B=\frac{\mu_0}{4\pi}\,
\frac{I\,\mathrm{d}\vec\ell\times\hat r}{r^{2}}.
\]
Para o \textbf{centro} da espira, todo elemento está à mesma distância $R$, e
$\mathrm{d}\vec\ell$ é sempre \emph{perpendicular} a $\hat r$ --- logo
$|\mathrm{d}\vec\ell\times\hat r|=\mathrm{d}\ell$.\\
(b) Além disso, todos os $\mathrm{d}\vec B$ apontam na \emph{mesma} direção
(perpendicular ao plano da espira), de modo que a integral vetorial vira
escalar:
\[
B=\frac{\mu_0I}{4\pi R^{2}}\oint\mathrm{d}\ell
=\frac{\mu_0I}{4\pi R^{2}}(2\pi R)
=\frac{\mu_0 I}{2R}.
\]
(c) \textbf{Três simplificações coincidiram:} distância constante, produto
vetorial de módulo máximo (ângulo reto) e contribuições todas paralelas. É a
combinação mais favorável possível.\\
\textbf{Fora do centro isso se perde:} no eixo, a distância deixa de ser $R$ e
as contribuições ganham componentes radiais que se cancelam --- exigindo
projeção. Fora do eixo, a integral não tem forma fechada elementar e recorre-se a
métodos numéricos ou a funções elípticas.}
\end{questao}

\begin{questao}[4]{}
\textbf{(Ampère contra Biot--Savart.)} Compare os dois métodos de cálculo de
campo magnético.
\begin{itensq}
\item Enuncie cada lei e sua condição de uso.
\item Determine quando a lei de Ampère é praticável.
\item Classifique quatro configurações quanto ao método adequado.
\end{itensq}
\resposta{Ampère exige simetria suficiente para tirar $B$ da integral;
Biot--Savart é geral mas trabalhoso}
\resolucao{(a) \textbf{Biot--Savart:}
$\mathrm{d}\vec B=\dfrac{\mu_0}{4\pi}
\dfrac{I\,\mathrm{d}\vec\ell\times\hat r}{r^{2}}$ --- vale \emph{sempre}, para
qualquer distribuição de corrente. É a lei fundamental.\\
\textbf{Ampère:} $\oint\vec B\cdot\mathrm{d}\vec\ell=\mu_0I_{env}$ --- também
sempre verdadeira, mas só \emph{útil} para calcular $B$ em casos especiais.\\
(b) \textbf{Ampère é praticável quando} existe uma curva fechada ao longo da
qual $|\vec B|$ é constante e $\vec B$ é paralelo (ou perpendicular) ao
caminho. Só assim $B$ pode ser retirado da integral:
\[
\oint\vec B\cdot\mathrm{d}\vec\ell=B\oint\mathrm{d}\ell=B\,\ell.
\]
Isso exige \textbf{simetria} --- cilíndrica, plana ou toroidal --- que precisa
ser \emph{conhecida de antemão}, por argumento físico.\\
(c)
\begin{center}
\begin{tabular}{lll}
\hline
Configuração & Método & Motivo\\
\hline
Fio retilíneo infinito & Ampère & simetria cilíndrica\\
Solenoide/toroide ideal & Ampère & simetria evidente\\
Espira circular (centro) & Biot--Savart & sem curva de $B$ constante\\
Segmento finito de fio & Biot--Savart & sem simetria alguma\\
\hline
\end{tabular}
\end{center}
\textbf{Erro comum:} tentar aplicar Ampère a uma espira circular usando a
própria espira como caminho. O campo \emph{não} é constante nem tangente ao
longo dela --- a integral existe, mas não permite isolar $B$.\\
\textbf{Analogia com a eletrostática:} Gauss está para Coulomb assim como
Ampère está para Biot--Savart. Em ambos os casos, a lei integral é elegante e
poderosa \emph{apenas} quando a simetria coopera.}
\end{questao}

\begin{questao}[4]{}
\textbf{(Projeto de solenoide.)} Projete um solenoide de $500$ espiras e
$\SI{25}{\centi\meter}$ para produzir $\SI{1}{\milli\tesla}$.
\begin{itensq}
\item Determine a corrente necessária.
\item Determine a potência dissipada, para fio de resistência total
      $\SI{2}{\ohm}$.
\item Avalie a uniformidade do campo.
\end{itensq}
\resposta{$I=\SI{0.398}{\ampere}$; $P=\SI{0.32}{\watt}$}
\resolucao{(a) $n=\dfrac{500}{0{,}25}=\SI{2000}{\per\meter}$, e
\[
I=\frac{B}{\mu_0 n}=\frac{\pot{1}{-3}}{(4\pi\times10^{-7})(2000)}
=\SI{0.398}{\ampere}.
\]
(b) $P=RI^{2}=2(0{,}398)^{2}=\SI{0.32}{\watt}$ --- desprezível, sem necessidade
de refrigeração.\\
(c) \textbf{Uniformidade.} Supondo diâmetro de $\SI{4}{\centi\meter}$, a razão
$L/D=6{,}25$ --- razoável, mas não excelente. Consequências:
\begin{itemize}
\item No \textbf{centro}, o campo fica cerca de $1$ a $2\%$ abaixo do valor
ideal $\mu_0nI$.
\item Nas \textbf{extremidades}, cai para aproximadamente metade.
\item A região com uniformidade melhor que $1\%$ limita-se ao terço central,
cerca de $\SI{8}{\centi\meter}$.
\end{itemize}
\textbf{Se for necessária melhor uniformidade:} alongar o solenoide (mantendo
$n$, o que exige mais espiras e mais tensão), ou adotar bobinas de Helmholtz.\\
\textbf{Verificação de coerência:} $\SI{1}{\milli\tesla}$ é vinte vezes o campo
terrestre --- o experimento precisa ser orientado ou compensado, sob pena de o
campo ambiente introduzir erro de $5\%$.}
\end{questao}

\begin{questao}[4]{}
\textbf{(Dipolo magnético.)} Uma espira de raio $R$ e corrente $I$ é observada
no eixo, a distância $x\gg R$.
\begin{itensq}
\item Obtenha a forma assintótica do campo.
\item Defina o momento de dipolo magnético e reescreva o resultado.
\item Compare com o dipolo elétrico e comente a ausência de monopolos.
\end{itensq}
\resposta{$B\to\dfrac{\mu_0 m}{2\pi x^{3}}$, com $m=I\pi R^{2}$}
\resolucao{(a) Partindo de
$B=\dfrac{\mu_0IR^{2}}{2(R^{2}+x^{2})^{3/2}}$ e fazendo $x\gg R$:
\[
B\to\frac{\mu_0IR^{2}}{2x^{3}}.
\]
(b) Definindo o \textbf{momento de dipolo magnético}
$m=I\,A=I\pi R^{2}$:
\[
B=\frac{\mu_0 m}{2\pi x^{3}}.
\]
\textbf{Verificação numérica:} para $R=\SI{10}{\centi\meter}$ e
$I=\SI{5}{\ampere}$, tem-se $m=\SI{0.157}{\ampere\meter\squared}$ e, a
$\SI{1}{\meter}$, $B=\pot{3.14}{-8}\,\si{\tesla}$. O valor exato é
$\pot{3.10}{-8}$ --- a aproximação erra $1{,}5\%$ a dez raios de distância.\\
(c) \textbf{Comparação com o dipolo elétrico.} Ambos decaem com $1/r^{3}$, e as
expressões são formalmente análogas ($p=2qa$ lá, $m=IA$ aqui).\\
\textbf{Mas há uma diferença fundamental.} O dipolo elétrico é feito de duas
cargas que \emph{podem} ser separadas --- e então cada uma produz campo
$1/r^{2}$. Já o dipolo magnético \textbf{não pode} ser separado: não existem
monopolos magnéticos.\\
\textbf{Consequência:} o dipolo é o termo \emph{dominante} de qualquer fonte
magnética estacionária. Não há termo $1/r^{2}$ em magnetostática --- e é por isso
que $\nabla\cdot\vec B=0$ figura entre as equações de Maxwell, sem análogo
elétrico.}
\end{questao}

\begin{questao}[4]{}
\textbf{(Blindagem magnética.)} Compare a blindagem de campos magnéticos com a
de campos elétricos.
\begin{itensq}
\item Explique por que uma gaiola metálica não blinda campo magnético estático.
\item Descreva os dois mecanismos disponíveis.
\item Discuta a dependência com a frequência.
\end{itensq}
\resposta{Não há "cargas magnéticas" a redistribuir; usa-se desvio por material
de alto $\mu_r$ ou correntes induzidas}
\resolucao{(a) \textbf{Por que a gaiola falha.} A blindagem elétrica funciona
porque o condutor possui \emph{cargas livres} que se redistribuem até anular o
campo interno. Não existe equivalente magnético: \textbf{não há monopolos
magnéticos livres} para migrar.\\
Um campo magnético \emph{estático} atravessa o cobre e o alumínio praticamente
sem atenuação.\\
(b) \textbf{Dois mecanismos disponíveis.}
\begin{itemize}
\item \textbf{Desvio de fluxo} (campos estáticos ou lentos). Um invólucro de
material de alto $\mu_r$ --- \emph{mu-metal}, com $\mu_r$ de dezenas de
milhares --- oferece caminho de baixa relutância. O fluxo prefere circular pelo
material e "contorna" a região protegida. Não há cancelamento: há \emph{desvio}.
\item \textbf{Correntes induzidas} (campos variáveis). Pela lei de Lenz, um
campo alternado induz correntes em um invólucro condutor, e essas correntes
geram campo oposto. A atenuação cresce com a frequência e com a condutividade.
\end{itemize}
(c) \textbf{Dependência com a frequência.}
\begin{center}
\begin{tabular}{lll}
\hline
Faixa & Mecanismo eficaz & Material\\
\hline
Estático (CC) & desvio de fluxo & mu-metal\\
Baixa (dezenas de Hz) & desvio, pouca indução & mu-metal\\
Alta (kHz e acima) & correntes induzidas & cobre, alumínio\\
\hline
\end{tabular}
\end{center}
\textbf{Consequência prática:} blindar campo elétrico é fácil e barato --- basta
qualquer metal aterrado. Blindar campo magnético estático é caro, exige material
especial, e a atenuação típica fica em uma ou duas ordens de grandeza, contra
muitas ordens no caso elétrico.\\
É por isso que salas de ressonância magnética têm blindagem magnética de
toneladas, enquanto a blindagem elétrica de um cabo é uma malha fina.}
\end{questao}

\begin{questao}[4]{}
\textbf{(Mapeamento experimental.)} Elabore um procedimento para mapear o campo
magnético de uma bobina usando sensor Hall.
\begin{itensq}
\item Descreva o procedimento e o tratamento do \emph{offset}.
\item Indique como obter as três componentes do campo.
\item Estabeleça o critério de validação dos dados.
\end{itensq}
\resposta{Medir com inversão de orientação para cancelar \emph{offset};
três orientações ortogonais para as componentes}
\resolucao{(a) \textbf{Procedimento.}
\begin{enumerate}
\item \textbf{Caracterizar o zero:} com a bobina \emph{desligada}, registrar a
leitura em cada ponto. Isso captura tanto o \emph{offset} do sensor quanto o
campo ambiente (terrestre e de equipamentos próximos).
\item \textbf{Medir com a bobina ligada} nos mesmos pontos.
\item \textbf{Subtrair} ponto a ponto. O resultado é o campo da bobina,
livre de \emph{offset} e de campo ambiente.
\end{enumerate}
\textbf{Alternativa mais robusta:} inverter o sentido da corrente na bobina e
tomar a \emph{semidiferença} das duas leituras. Isso cancela tudo o que não
depende da corrente.\\
(b) \textbf{Três componentes.} O sensor Hall responde apenas à componente
perpendicular à sua face. Para obter $\vec B$ completo, meça em três orientações
ortogonais no mesmo ponto:
\[
|\vec B|=\sqrt{B_x^{2}+B_y^{2}+B_z^{2}}.
\]
Um sensor triaxial faz isso simultaneamente e evita erros de reposicionamento.\\
(c) \textbf{Critérios de validação.}
\begin{itemize}
\item \textbf{Linearidade com a corrente:} $B$ deve ser proporcional a $I$.
Meça em três correntes e verifique --- desvio indica saturação de núcleo ou
não linearidade do sensor.
\item \textbf{Simetria:} pontos simétricos em relação ao eixo devem dar valores
iguais. Discrepância revela desalinhamento do sensor ou da bobina.
\item \textbf{Valor central:} deve concordar com $\mu_0nI$ dentro da incerteza,
descontada a correção de comprimento finito.
\item \textbf{Reprodutibilidade:} repita o primeiro ponto ao final. Diferença
indica deriva térmica do sensor.
\end{itemize}
\textbf{Registre sempre a temperatura} --- a sensibilidade Hall varia
tipicamente algumas centenas de ppm por grau.}
\end{questao}

\begin{questao}[4]{}
\textbf{(Limites de eletroímãs.)} Analise o que limita, na prática, o campo
magnético alcançável.
\begin{itensq}
\item Determine o limite imposto pela saturação do núcleo.
\item Determine o limite térmico de bobinas de cobre.
\item Descreva as soluções empregadas para superá-los.
\end{itensq}
\resposta{Núcleos saturam entre $1{,}5$ e $\SI{2}{\tesla}$; acima disso é
preciso bobina sem núcleo, com corrente muito alta}
\resolucao{(a) \textbf{Saturação.} Materiais ferromagnéticos alinham seus
domínios completamente em torno de $1{,}5$ a $\SI{2}{\tesla}$ (ferro-silício,
$\approx\SI{2}{\tesla}$; ferrites, bem menos). Acima disso, $\mu_r$ cai para
próximo de $1$ e o núcleo deixa de contribuir --- \textbf{acrescentar corrente
passa a render tão pouco quanto no ar}.\\
Este é um limite \emph{de material}, não de projeto: nenhuma geometria o
contorna.\\
(b) \textbf{Limite térmico.} Sem núcleo, $B=\mu_0nI$ exige corrente enorme. A
potência dissipada vai com $I^{2}$, e a densidade de corrente admissível em
cobre refrigerado a água fica em torno de
$\SI{20}{\ampere\per\milli\meter\squared}$.\\
Eletroímãs resistivos de laboratório atingem $1$ a $\SI{2}{\tesla}$ consumindo
\emph{megawatts} --- e boa parte do custo operacional é a refrigeração.\\
(c) \textbf{Soluções empregadas.}
\begin{itemize}
\item \textbf{Supercondutores:} resistência nula elimina o limite térmico.
Bobinas de NbTi e $\mathrm{Nb_3Sn}$ chegam a $\SI{20}{\tesla}$ e são a base de
equipamentos de ressonância magnética e de aceleradores. O custo passa a ser a
criogenia.
\item \textbf{Campos pulsados:} correntes altíssimas por milissegundos, antes
que o calor se acumule. Alcançam $\SI{100}{\tesla}$, com a bobina
frequentemente destruída no processo.
\item \textbf{Ímãs híbridos:} bobina supercondutora externa somada a uma
resistiva interna --- combinação que detém os recordes de campo contínuo, acima
de $\SI{45}{\tesla}$.
\end{itemize}
\textbf{Limite último:} acima de algumas dezenas de tesla, a \emph{pressão
magnética} $B^{2}/2\mu_0$ supera a resistência mecânica de qualquer material. A
$\SI{100}{\tesla}$ ela vale cerca de $\SI{4}{\giga\pascal}$ --- a bobina
literalmente se rompe.}
\end{questao}

\begin{questao}[4]{}
\textbf{(Superposição e simetria.)} Um condutor cilíndrico de raio $a$ possui
uma cavidade cilíndrica de raio $b$, paralela ao eixo e deslocada de $d$.
A corrente $I$ distribui-se uniformemente na seção restante.
\begin{itensq}
\item Formule a estratégia de solução por superposição.
\item Determine o campo dentro da cavidade.
\item Interprete o resultado.
\end{itensq}
\resposta{O campo na cavidade é \textbf{uniforme}, de módulo
$\mu_0Jd/2$}
\resolucao{(a) \textbf{Estratégia.} O condutor furado equivale à
\emph{superposição} de:
\begin{itemize}
\item um cilindro \textbf{maciço} de raio $a$ com densidade de corrente
$+J$; \emph{mais}
\item um cilindro de raio $b$, na posição da cavidade, com densidade
$-J$.
\end{itemize}
A soma reproduz corrente $J$ na região sólida e zero na cavidade --- exatamente
a situação real.\\
(b) Dentro de um cilindro maciço de densidade $J$, a lei de Ampère dá
$B=\dfrac{\mu_0 J r}{2}$, ou vetorialmente
$\vec B=\dfrac{\mu_0}{2}\,\vec J\times\vec r$, com $\vec r$ medido a partir do
\emph{eixo daquele} cilindro.\\
Somando as duas contribuições em um ponto da cavidade, com $\vec r_1$ e
$\vec r_2$ medidos a partir de cada eixo:
\[
\vec B=\frac{\mu_0}{2}\vec J\times\vec r_1
-\frac{\mu_0}{2}\vec J\times\vec r_2
=\frac{\mu_0}{2}\vec J\times(\vec r_1-\vec r_2)
=\frac{\mu_0}{2}\vec J\times\vec d.
\]
\textbf{O vetor $\vec d$ é constante} --- ele liga os dois eixos e não depende do
ponto.\\
(c) \textbf{Interpretação: o campo na cavidade é rigorosamente uniforme}, de
módulo $\mu_0Jd/2$ e direção perpendicular a $\vec d$.\\
Isso é notável: uma geometria sem simetria evidente produz, na região vazia, o
campo mais simples possível.\\
\textbf{Aplicação:} o mesmo resultado, com $\vec J$ substituído por densidade de
carga, dá campo elétrico uniforme na cavidade de uma esfera carregada --- e é
usado no projeto de aceleradores e de ímãs de deflexão, em que se busca campo
uniforme em uma região de acesso.}
\end{questao}
